\documentclass[10pt]{article}
\usepackage[T1]{fontenc}
\usepackage[utf8]{inputenc}
\usepackage[french]{babel}
\usepackage{amsmath,amssymb,mathrsfs,mathpazo}
\usepackage{geometry}
\geometry{a4paper,margin=2cm}
\usepackage{graphicx,array,tabularx,tabularray,colortbl,multirow,multicol,booktabs,enumitem,pifont,ulem,fancybox,diagbox,cancel}
\usepackage{pstricks,pstricks-add,pst-plot,pst-tree,pst-node,pst-text,pst-func,pst-math,pst-bspline,pst-3dplot,pst-eucl,pst-solides3d}
\let\pscalculate\undefined
\usepackage{pst-bezier}
\usepackage{tikz,pgfplots,tkz-tab}
\usetikzlibrary{arrows,arrows.meta,patterns,calc,positioning,trees,decorations.markings,shapes.geometric,shapes.misc,shapes.arrows,intersections,decorations.pathmorphing,decorations.pathreplacing,backgrounds,shadows,babel}
\pgfplotsset{compat=1.18}
\usepgfplotslibrary{fillbetween}
\usepackage[np]{numprint}
\usepackage{listings,makecell,pas-tableur,twoopt}
\let\barre\undefined
\usepackage[pstricks]{tablvar}
\definecolor{darkgreen}{RGB}{0,100,0}
\providecommand{\red}{\color{red}}
\providecommand{\blue}{\color{blue}}
\providecommand{\green}{\color{green}}
\providecommand{\black}{\color{black}}
\providecommand{\white}{\color{white}}
\providecommand{\gray}{\color{gray}}
\providecommand{\cyan}{\color{cyan}}
\providecommand{\magenta}{\color{magenta}}
\providecommand{\yellow}{\color{yellow}}
\providecommand{\violet}{\color{violet}}
\providecommand{\orange}{\color{orange}}
\providecommand{\euro}{\texteuro}
\providecommand{\EUR}{\texteuro}
\providecommand{\ieme}{\textsuperscript{e}}
\providecommand{\degres}{\ensuremath{^\circ}}
\providecommand{\degree}{\ensuremath{^\circ}}
\providecommand{\vv}[1]{\overrightarrow{#1}}
\providecommand{\Coord}[1]{#1}
\providecommand{\decosix}{}
\providecommand{\starredbullet}{\ensuremath{\bullet}}
\pagestyle{empty}
\setlength{\parindent}{0pt}
\usepackage{adjustbox}\DeclareUnicodeCharacter{2236}{:}\DeclareUnicodeCharacter{2190}{\ensuremath{\leftarrow}}
\providecommand{\R}{\mathbb{R}}
\providecommand{\N}{\mathbb{N}}
\providecommand{\Z}{\mathbb{Z}}
\providecommand{\Q}{\mathbb{Q}}
\providecommand{\D}{\mathbb{D}}
\providecommand{\C}{\mathbb{C}}
\providecommand{\e}{{\mathrm{e}}}
\providecommand{\vect}[1]{\overrightarrow{#1}}
\providecommand{\vectt}[1]{\overrightarrow{\mathrm{#1}}}
\providecommand{\vvt}[1]{\overrightarrow{\mathrm{#1}}}
\providecommand{\np}[2][]{{#2~#1}}
\providecommand{\vv}[1]{{\overrightarrow{#1}}}
\renewcommand{\d}{{\mathrm{d}}}
\providecommand{\degres}{{^\circ}}
\providecommand{\degre}{{^\circ}}
\providecommand{\blue}{\color{blue}}
\providecommand{\red}{\color{red}}
\providecommand{\enskip}{\quad}
\providecommand{\llbracket}{\lbrack\!\lbrack}
\providecommand{\rrbracket}{\rbrack\!\rbrack}
\providecommand{\euro}{{\text{€}}}
\providecommand{\no}{{\text{n°}}}
\providecommand{\interval}[2]{{[#1;#2]}}
\providecommand{\intervFO}[2]{{[#1;#2[}}
\providecommand{\dotfill}{\ldots}

\providecommand{\ds}{\displaystyle}
\providecommand{\textding}[1]{#1}
\providecommand{\Oij}{\(\mathrm{O};\vec{\imath},\vec{\jmath}\)}
\providecommand{\Oijk}{\(\mathrm{O};\vec{\imath},\vec{\jmath},\vec{k}\)}
\providecommand{\Ouv}{\(\mathrm{O};\vec{u},\vec{v}\)}
\providecommand{\intervFF}[2]{\ensuremath{[#1;#2]}}
\renewcommand{\euro}{\ifmmode\text{\texteuro}\else\texteuro\fi}
\DeclareUnicodeCharacter{00B0}{\ifmmode^{\circ}\else\textdegree\fi}
\DeclareUnicodeCharacter{20AC}{\ifmmode\text{\texteuro}\else\texteuro\fi}
% Énoncés : Éducation nationale. Transcription et corrigés : APMEP, Denis Vergès et contributeurs. Découpage et classement : CoopMaths, Sébastien Lozano, Éric Elter et contributeurs. Adaptation de la présentation : Maths Explorer.
% CC BY-SA : https://coopmaths.fr/www/about/
% https://coopmaths.fr/alea/static/bac/2024/tex/bac_2024_06_sujet1_metropole_3_cor.tex
\begin{document}
\noindent\textit{APMEP / CoopMaths. CC BY-SA. Présentation adaptée par Maths Explorer.}\par\medskip
\medskip

\begin{enumerate}
\item
	\begin{enumerate}
		\item On a $\vect{\text{CA}}\begin{pmatrix}5\\5\\-10\end{pmatrix}$ et $\vect{\text{DC}}\begin{pmatrix}0\\0\\\frac{25}{2}\end{pmatrix}$.
		
Or ces deux vecteurs ne sont pas colinéaires et :

$\vect{n_1} \cdot \vect{\text{CA}} = 5 - 5 + 0 = 0$ et 

$\vect{n_1} \cdot \vect{\text{DC}} = 0 +  0 + 0 = 0$.

Conclusion : le vecteur $\vect{n_1}$ est orthogonal à deux vecteurs non colinéaires du plan (CAD) : il est donc normal à ce plan.
		\item On sait qu'alors :
		
$M(x~;~y~;~z) \in (\text{CAD}) \iff 1x - 1y + 0z + d = 0 \iff x - y + d = 0$, avec \:$ d \in \R$.

Ceci est vrai pour exemple pour C :

$C(0~;~0~;~10) \in (\text{CAD}) \iff 0 - 1y + 0 + d = 0$  avec \: $d \in \R$, soit $d = 0$.

Donc $M(x~;~y~;~z) \in (\text{CAD}) \iff x - y  = 0$.
	\end{enumerate}
\item
	\begin{enumerate}
		\item Si la droite $(D)$ est sécante au plan (CAD) en un point H les coordonnées de ce point vérifient les équations paramétriques de ($D$) et celle du plan soit :

H$(x~;~y~;~z) \in D \cap (\text{CAD})\iff \left\{\begin{array}{l c l}
x&=&\phantom{5 -} + \frac52 t\\
y &=&5 - \frac52 t\\
z&=&\phantom{5 -}0t\\
x - y = 0
\end{array}\right.\:(t \in \R) \Rightarrow \dfrac52t - \left(5 - \dfrac52 t \right) = 0 \iff 5t - 5 = 0 \iff t = 1$.

Donc H$\left(\dfrac52~;~\dfrac52~;~0\right)$.
		\item On a $M(x~;~y~;~z) \in D \iff \left\{\begin{array}{l c l}
x&=&0 + \frac52 t\\
y &=&5 - \frac52 t\\
z&=&0 + 0t
\end{array}\right.$

On reconnait dans les coefficients de $t$ les coordonnées du vecteur $\begin{pmatrix}\frac52\\- \frac52\\0\end{pmatrix}$, soit $\frac52\vect{n_1}$ et (0~;~5~;~0) sont les coordonnées de B.

Donc $D$ est la droite contenant B et de vecteur directeur $\vect{n_1}$.

Or on a vu que le vecteur $\vect{n_1}$ est normal au plan (CAD), donc la droite $D$ est perpendiculaire à (CAD) au point H.
	\end{enumerate}
\item 
	\begin{enumerate}
		\item On a vu que (BH) = $D$ donc a pour vecteur directeur $\vect{n_1}$.
		
$\vect{\text{HA}}\begin{pmatrix}\frac52\\\frac52\\0\end{pmatrix}$ et $\vect{n_1} \cdot \vect{\text{HA}} = \dfrac52 - \dfrac52 + 0 = 0$.
		
Les vecteurs $\vect{\text{BH}}$ et $\vect{\text{HA}}$ sont orthogonaux donc les droites (BH) et (HA) sont perpendiculaires, le triangle ABH est rectangle en H.
		\item On a $\left\|\vect{\text{BH}}\right\|^2 = \frac{25}{4} + \frac{25}{4} + 0 = \frac{50}{4}$, d'où $\left\|\vect{\text{BH}}\right\| = \text{BH} = \sqrt{\frac{50}{4}} = \dfrac{5\sqrt 2}{2}$ et 
		
$\left|\vect{\text{HA}}\right\|^2 = \frac{25}{4} + \frac{25}{4} + 0 = \frac{50}{4}$, d'où $\left\|\vect{\text{HA}}\right\| = \text{HA} = \sqrt{\frac{50}{4}} = \dfrac{5\sqrt 2}{2}$.

Donc $\mathcal{A}(\text{ABH}) = \dfrac{\text{BH} \times \text{HA}}{2} = \dfrac12 \times \dfrac{5\sqrt 2}{2} \times \dfrac{5\sqrt 2}{2} = \dfrac{25 \times 2}{8} = \dfrac{25}{4}$.
	\end{enumerate}
\item 
	\begin{enumerate}
		\item Les points A, B et H ont une cote nulle : ils appartiennent donc au plan définis par les vecteurs unitaires $\vect{\imath}$ et $\vect{\jmath}$ et C a une abscisse et une ordonnée nulles, donc appartient à l'axe $\left(\text{O},~\vect{k}\right)$ qui est perpendiculaire au plan (ABH), donc (CO) est la hauteur contenant C du tétraèdre ABCH.
		\item De façon évidente CO $ = 10$, donc :
		
		$\mathcal{V}(\text{ABCH}) = \dfrac{\mathcal{A}(\text{ABH}) \times \text{CO}}{3} = \dfrac{\frac{25}{4} \times 10}{3} = \dfrac{125}{6}$.
	\end{enumerate}
\item On a $\vect{\text{BA}}\begin{pmatrix}5\\0\\0\end{pmatrix}$ et $\vect{\text{BC}}\begin{pmatrix}0\\-5\\10\end{pmatrix}$, d'où $\vect{\text{BA}} \cdot \vect{\text{BC}} = 0 + 0 + 0 = 0$ : les vecteurs sont orthogonaux donc les droites (BA) et (BC) sont perpendiculaires, le triangle (ABC) est rectangle en B.

On a BA$^2 = 5^2$, donc BA $ = 5$ et BC$^2 = 5^2 + 10^2 = 25 + 100 = 125$, donc BC $= \sqrt{125} = 5 \sqrt 5$.

L'aire du triangle (ABC) est égale à :

$\mathcal{A}(\text{ABC}) = \dfrac{\text{AB} \times \text{BC}}{2} = \dfrac{ 5 \times 5 \sqrt 5}{2} = \dfrac{25\sqrt 5}{2}$.

On a $\mathcal{V}(\text{ABCH}) = \dfrac{\mathcal{A}(\text{ABC}) \times d}{3}$, $d$ étant la distance du point H au plan (ABC).

On a donc : $\dfrac{125}{6} = \dfrac{\dfrac{25\sqrt 5}{2} \times d}{3} \iff d = \dfrac{125}{6} \times \dfrac{3 \times 2}{25\sqrt 5} = \sqrt 5$.
\end{enumerate}

\bigskip

\end{document}
